%% P23 \dash{} Random variables and distributions
%%
%% Stub. The brief below is the contract for this program: what it must argue,
%% and what it must buy the reader. Writing the program means deleting the
%% \programstub{} block and replacing it with the program. If the brief turns
%% out to be wrong once the mathematics has been worked, change it and record
%% the contradiction in CLAUDE.md under *Resolved questions*.

\program{Random variables and distributions}
\label{prog:P23}

\programstub{%
Argument: a random variable is a function on the sample space; expectation
and variance are its first two summaries; six distributions cover almost
everything in this field. Payoff: sampling a token is a draw from a
categorical distribution, and temperature, top-k and top-p are three ways of
editing that distribution before drawing \dash{} described as distribution
surgery, with what each destroys; the Gumbel-max trick, so the reader can
see that \code{argmax(logits + gumbel)} is exact categorical sampling and
not an approximation; the Gaussian introduced through its role rather than
its formula; expectation as the linear operator that makes almost every
derivation later in the book short.

\medskip
\textbf{Planned length:} 60 frames.
}
